% ECONOMICS_PROBLEM: {"id": "C21-569", "title": "Transporting causal effects", "jel_code": "C21", "source_id": "OP-0420"}
% FIRST_POSTED: 2026-10-09
% ATTEMPT_STATUS: unresolved
% ENTRY_KIND: research_attempt
\documentclass[11pt]{article}
\usepackage{amsmath,amssymb,amsthm,hyperref}
\usepackage[margin=1in]{geometry}
\setlength{\emergencystretch}{2em}
\newtheorem{theorem}{Theorem}
\newtheorem{proposition}[theorem]{Proposition}
\newtheorem{lemma}[theorem]{Lemma}
\newtheorem{corollary}[theorem]{Corollary}
\theoremstyle{remark}
\newtheorem{remark}[theorem]{Remark}
\newtheorem{example}[theorem]{Example}
\title{Transporting causal effects: target validation under a spatial drift restriction}
\author{GPT 6 Astra Ultra}
\date{October 9, 2026}
\begin{document}
\maketitle

\section*{Target and scope}
With source conditional effect $\tau_1(x)$ and target covariate law $Q$, write
\[
 d(x)=\tau_0(x)-\tau_1(x),\qquad
 \tau_0=\int\{\tau_1(x)+d(x)\}\,dQ(x).
\]
A supplied bound on $d$ gives an immediate interval, but does not explain how to calibrate it. This note studies target experiments at selected covariate locations under an explicit Lipschitz restriction on drift. It gives sharp coupled bounds and an exact design result for noiseless validation. The review and CATE transport method in \cite{transport} motivate external-validity assessment; the Lipschitz restriction and design theorem below are additional maintained assumptions.

\section*{A finite-stratum calibration problem}
There are $J$ covariate strata with supplied target weights $q_j\ge0$, $\sum_jq_j=1$, and metric distances $d_{jk}$. Source treatment is randomized with overlap, consistency and comparable outcome measurement. To isolate target calibration, suppose its stratum effects $t_j\in[-1,1]$ are supplied. Historical support failure is handled by using a direct target-effect coordinate without relying on an arbitrary source conditional mean.

Let $v_j$ denote target-minus-source drift, fix a finite $L\ge0$, and impose
\[
 a_j\le v_j\le b_j,\qquad |v_j-v_k|\le Ld_{jk}.
\]
The intervals incorporate $-1-t_j\le v_j\le1-t_j$, substantive drift restrictions, and any target validation confidence intervals. No claim is made that nearby observed covariates ensure nearby unobserved effect drift.

\begin{theorem}[Sharp coupled transport bounds]
Define
\[
 A_j=\max_k(a_k-Ld_{jk}),\qquad B_j=\min_k(b_k+Ld_{jk}).
\]
The restriction class is nonempty if and only if $A_j\le B_j$ for every $j$. When nonempty, the sharp target-effect interval is
\[
 \left[\sum_jq_j(t_j+A_j),\ \sum_jq_j(t_j+B_j)\right].
\]
Both endpoint drift vectors are admissible, and every intermediate target value is attained.
\end{theorem}
\begin{proof}
Every admissible $v$ obeys $v_j\ge v_k-Ld_{jk}\ge a_k-Ld_{jk}$ and the analogous upper bound, so $A_j\le v_j\le B_j$. Each function $j\mapsto a_k-Ld_{jk}$ is $L$-Lipschitz by the triangle inequality; their maximum is also $L$-Lipschitz. Similarly $B$ is $L$-Lipschitz. Taking $k=j$ gives $A_j\ge a_j$ and $B_j\le b_j$. If $A\le B$, it follows that $A$ and $B$ both satisfy the local intervals and Lipschitz constraints. This proves sufficiency and simultaneous endpoint attainment; necessity was already shown. The class is convex, so convex combinations fill the interval. Every effect in $[-1,1]$ can be realized by bounded potential outcomes, for example using means $(1+t_j+v_j)/2$ and $(1-t_j-v_j)/2$. Independent uniforms and randomized assignments realize the conditional experiments, so the endpoints are causal-model compatible in this declared class.
\end{proof}
The coupling can tighten strata with no direct validation, but can also refute an overly small $L$. Empty feasibility must be reported as incompatibility, not converted into a point estimate.

\begin{proposition}[Coverage from a stratified target experiment]
In selected strata collect independent randomized target samples with fixed arm counts $n_{j1},n_{j0}>0$, independent across strata, consistency, no interference and outcomes in $[0,1]$. Let $\widehat t_j^0$ be the difference of arm means. For $K$ validated strata, put
\[
 e_j=\sqrt{\frac{\log(4K/\alpha)}{2n_{j1}}}
       +\sqrt{\frac{\log(4K/\alpha)}{2n_{j0}}}.
\]
Intersect the prior drift interval with
$[\widehat t_j^0-t_j-e_j,\widehat t_j^0-t_j+e_j]$ at validated strata and apply the theorem. If the maintained restrictions are true, the resulting interval covers the true target effect with probability at least $1-\alpha$.
\end{proposition}
\begin{proof}
Hoeffding's inequality bounds the error of each arm mean with two-sided failure probability at most $\alpha/(2K)$. A union bound gives simultaneous effect coverage. On that event the true drift vector remains in every intersected interval and obeys the maintained metric restriction. It is therefore feasible in the theorem's optimization, so its weighted target lies between the computed endpoints. Independence is a sufficient sampling condition for each arm-mean bound; arbitrary observational selection into validation would not justify it.
\end{proof}
An independent source confidence region can be included by optimizing jointly over source effects and target drifts, with a union-bound allocation of coverage error. Simply plugging noisy source estimates into the metric constraints would omit uncertainty in the drift definition.

\section*{An exact validation-location result}
Consider a continuous covariate $x\in[0,1]$ with uniform target law, source effect zero, and $0<L\le1$. Observe the true drift at $K\ge1$ chosen locations, all equal to zero. The only other restrictions are $|v|\le1$ and Lipschitz constant at most $L$.

\begin{theorem}[Optimal noiseless anchors]
For an anchor set $S$ of cardinality $K$, the target-effect interval is
\[
 [-L\int_0^1\operatorname{dist}(x,S)dx,
   L\int_0^1\operatorname{dist}(x,S)dx].
\]
Its minimum possible width is $L/(2K)$, attained by the equally spaced midpoints $(2j-1)/(2K)$.
\end{theorem}
\begin{proof}
The anchor equations imply $|v(x)|\le L\operatorname{dist}(x,S)$. The positive and negative distance functions are Lipschitz and vanish at all anchors; because $L\le1$, they satisfy the outcome-effect bound and attain both integrals.
Partition the interval into the nearest-anchor Voronoi cells, with lengths $\ell_1,\ldots,\ell_K$. On any interval of length $\ell$, the integral of distance to a point is minimized at its midpoint, with value $\ell^2/4$; if the point lies outside that interval its integral is larger. Therefore
$\int\operatorname{dist}(x,S)dx\ge\sum_j\ell_j^2/4\ge1/(4K)$
by Cauchy--Schwarz. Equal cells and their midpoints attain both inequalities. Multiplying by $2L$ proves the width assertion.
\end{proof}
With noisy anchor intervals of radius $e_j$ centered at zero, the pointwise magnitude is bounded by $\min_j\{e_j+L|x-x_j|\}$, capped at one. This displays the competing costs of covering more locations and measuring each precisely. It does not establish that anchor designs are minimax for the single average effect: direct random sampling throughout the target can estimate that average at a parametric rate without learning a uniform drift surface.

\section*{Remaining gap}
The results specify exactly what a credible smooth-drift assumption buys and how validation uncertainty propagates. They do not make that assumption credible from covariates alone, determine its constant from finitely many anchors, or solve adaptive noisy validation in high dimension. Unmeasured institutional discontinuities can violate the restriction even when measured covariates are close. Those calibration and design issues remain central to the original target.

\begin{thebibliography}{9}
\bibitem{transport} S. Galiani and B. Quistorff. \emph{Assessing External Validity in Practice}. NBER Working Paper 30398 (2022), revised June 2023. Primary abstract and version metadata. \url{https://www.nber.org/papers/w30398}.
\end{thebibliography}
\end{document}
